\documentclass{rapportECL}
\usepackage{lipsum}

\begin{document}

%----------- Informations du rapport ---------

\titre{Art of Data Science} %Titre du fichier .pdf


%----------- Initialisation -------------------
        
%\fairemarges %Afficher les marges
\fairepagedegarde %Créer la page de garde
\tabledematieres %Créer la table de matières
\tableofcontents
\newpage
%------------ Corps du rapport ----------------








\section{Introduction}

The way organizations make decisions has been evolving for half a century. Before the introduction of Business Intelligence, the only options were gut instinct, loudest voice, and best argument. Sadly, this method still exists today, and in some pockets it is the predominant means by which the organization acts. Take our advice and never, ever work for such a company!

Fortunately for our economy, most organizations began to inform their decisions with real information through the application of simple statistics. Those that did it well were rewarded; those that did not failed. We are outgrowing the ability of simple stats to keep pace with market demands, however. The rapid expansion of available data, and the tools to access and make use of the data at scale, are enabling fundamental changes to the way organizations make decisions.

Data Science is required to maintain competitiveness in the increasingly data-rich environment. Much like the application of simple statistics, organizations that embrace Data Science will be rewarded while those that do not will be challenged to keep pace. As more complex, disparate data sets become available, the chasm between these groups will only continue to widen. The figure,
The Business Impacts of Data Science, highlights the value awaiting organizations that embrace Data Science.

As a result of the increasingly interconnected world, huge amounts of data are being generated and stored every instant. Data Science can be used to transform data into insights that help improve existing processes. Operating costs
can be driven down dramatically by effectively incorporating the complex interrelationships in data like never before. This results in better quality assurance, higher product yield and more effective operations.


\textbf{We don't want to make decisions entirely on gut feelings, rather we want to bring in evidence in the form of data and insights. We want to combine them with domain expertise and make informed decisions}

\section{Defining Data Science}

Data Science is the art of turning data into actions. This is accomplished through the creation of data products, which provide actionable information without exposing decision makers to the underlying data or analytics (e.g., buy/sell strategies for financial instruments, a set of actions to improve product yield, or steps to improve product marketing).

Performing Data Science requires the extraction of timely, actionable information from diverse data sources to drive data products. Examples of data products include answers to questions such as: “Which of my products should I advertise more heavily to increase profit? How can I improve my compliance program, while reducing costs? What manufacturing process change will allow me to build a better product?” The key to answering these questions is: understand the data you have and what the data inductively tells you

In short,  Data Science is the art of turning data into actions. process, tools and technologies for humans and computers to work together to transform data into insights. Data products provide actionable information without exposing decision makers to the underlying data or analytics (e.g., buy/sell strategies for financial instruments, a set of actions to improve product yield, or steps to improve product marketing).




\section{What makes Data Science Different?}

Data Science supports and encourages shifting between deductive (hypothesis-based) and inductive (pattern-based) reasoning. This is a fundamental change from traditional analytic approaches. Inductive reasoning and exploratory data analysis provide a means to form or refine hypotheses and discover new analytic paths. In fact, to do the discovery of significant insights that are the hallmark of Data Science, you must have the tradecraft and the interplay between inductive
and deductive reasoning. By actively combining the ability to reason deductively and inductively, Data Science creates an environment where models of reality no longer need to be static and empirically based. Instead, they are constantly tested, updated and improved until better models are found. These concepts are summarized in the figure, The Types of Reason and Their Role in Data Science Tradecraft.

Data Science supports and encourages shifting between deductive (hypothesis-based) and inductive (pattern- based) reasoning. Inductive reasoning and exploratory data analysis provide a means to form or refine hypotheses and discover new analytic paths. Models of reality no longer need to be static. They are constantly tested, updated and improved until better models are found.

The differences between Data Science and traditional analytic approaches do not end at seamless shifting between deductive and inductive reasoning. Data Science offers a distinctly different perspective than capabilities such as Business Intelligence. However, data science should not replace Business Intelligence functions within
an organization. The two capabilities are additive and complementary, each offering a necessary view of business operations and the operating environment. 

The figure, Business Intelligence and Data Science – A Comparison, highlights the differences between the two capabilities. 


Key contrasts include:

\begin{itemize}
    \item Discovery vs. Pre-canned Questions: Data Science actually works on discovering the question to ask as opposed to just asking it.
    \item Power of Many vs. Ability of One: An entire team provides a common forum for pulling together computer science, mathematics and domain expertise.
    \item  Prospective vs. Retrospective: Data Science is focused on obtaining actionable information from data as opposed to reporting historical facts.
\end{itemize}


\textbf{what is different from traditional reporting is this idea that we have sometimes have ambiguty around what question to ask? So it requires sometimes exploratory data analysism talking to domain experts, to first forumulate the right questions. and that is sometimes the most important bit.} 

\textbf{And then another thing is the fractal nature of the whole process. This means one question might lead us to ask more questions or we might go back and reformulate questions in a better way.}


\section{The Data Science Process}

The Four Key Activities of a Data Science Endeavor. Data Science purists will likely disagree with this approach


\begin{itemize}
    
    
    \item \textbf{Activity 1: Acquire:} This activity focuses on obtaining the data you need. Given the nature of data, the details of this activity depend heavily on who you are and what you do. 

    \item \textbf{Activity 2: Prepare} Great outcomes don’t just happen by themselves. A lot depends on preparation, and in Data Science, that means manipulating the data to fit your analytic needs. This stage can consume a great deal of time, but it is an excellent investment. The benefits are immediate and long term.

    \item \textbf{Activity 3: Analyze} This is the activity that consumes the lion’s share of the team’s attention. It is also the most challenging and exciting (you will see a lot of ‘aha moments’ occur in this space). 

    \item \textbf{Activity 4: Act} Every effective Data Science team analyzes its data with a purpose – that is, to turn data into actions. Actionable and impactful insights are the holy grail of Data Science. Converting insights into action can be a politically charged activity, however.
    
\end{itemize}





\section{Journey of Developing Capabilities}

Data Science capabilities can be built over time. Organizations mature through a series of stages – Collect, Describe, Discover, Predict, Advise – as they move from data deluge to full Data Science maturity. At each stage, they can tackle increasingly complex analytic goals with a wider breadth of analytic capabilities. However, organizations need not reach maximum Data Science maturity to achieve success. Significant gains can be found in every stage.


Data Science is a different kind of team sport. Data Science teams need a broad view of the organization. Leaders must be key advocates who meet with stakeholders to ferret out the hardest challenges, locate the data, connect disparate parts of the business, and gain widespread buy-in.





\section{Classes of Analytic Techniques}

\subsection{Describe}


\subsection{Discover}


\subsection{Predict}


\subsection{Advice}


\section{Fractal Analytics Model}









\section{Case Studies}


Book: Data Science on AWS

Mackinsey Analytics https://www.databricks.com/resources/ebook/the-big-book-of-data-science-use-cases-nurture 

https://www.databricks.com/resources/ebook/big-book-of-machine-learning-use-cases 

\section{Analytics Patterns}


1. Aggregation
2. Experimentation
3. Prediction
4. Clustering
5.. Decision trees.
6. Cumulation and derivatives
7. Funnel analysis.


\section{Resources}

1. Data Science: The Hard Parts: Techniques for Excelling at Data Science

2. https://jtleek.com/ads2020/ 

3. Art of Data Science book 

4. Field Guide etc. etc. 

\end{document}
